\documentclass[11pt,a4paper]{appletmanual}
\usepackage[spanish,es-noquoting,es-nodecimaldot]{babel}
\manuallang{es}

\appletname{Impuesto a la Renta del Trabajo}
\appletsub{Un impuesto por tramos dobla la restricción entre consumo y ocio.
Con tipos marginales la frontera se quiebra; con tipos medios se rompe, y
cruzar un umbral deja al trabajador estrictamente peor.}
\appletdir{labor-tax}

\begin{document}
\manualtitlepage{Manual del applet · Microeconomía I · Universidad de los Andes\\
Traducido del applet de GeoGebra «Restricción Presupuestaria con Impuesto a la Renta».\\
Código MIT. Se abre en el navegador; no hay nada que instalar.}

\manualtoc
\thispagestyle{empty}
\clearpage
\setcounter{page}{1}

% ==========================================================================
\section{Qué es esto}

Un trabajador reparte una dotación de tiempo entre trabajar y no trabajar. Si
dedica la fracción $L$ a trabajar, gana $wL$, paga un impuesto sobre esa renta
y se gasta lo que queda en consumo:
\[
  c(L) \;=\; \frac{m_0 + wL - T(wL)}{p_x},
  \qquad \ell = 1 - L .
\]

La dotación está normalizada a $1$, así que $L$ es la fracción del tiempo que
se trabaja y $\ell$ el ocio. Lo único que hay que elegir es $L$.

Todo el interés está en $T(\cdot)$, el impuesto por tramos, y en una decisión de
diseño que los sistemas fiscales reales toman y que casi nunca se discute en
clase: si el tipo de cada tramo se aplica \emph{solo a la renta que cae dentro
de ese tramo} o \emph{a toda la renta}. Los dos esquemas suenan casi igual y
producen fronteras presupuestarias completamente distintas.

\figher[1.0]{marginal}{Tipos marginales. El impuesto de cada tramo grava solo
la renta que cae dentro de él. La renta neta es continua, así que la frontera
se \emph{quiebra} pero no se corta.}

\begin{amkey}
Con tipos marginales la frontera tiene un \textbf{quiebre}: la pendiente cambia
de golpe pero la línea no se rompe. Con tipos medios tiene un
\textbf{salto}: ganar una unidad más deja al trabajador con menos. Un quiebre
atrae óptimos; un salto los empuja hacia abajo y los deja pegados justo por
debajo del umbral.
\end{amkey}

\subsection{Para quién}

Para el tema de elección entre consumo y ocio, o para una clase de economía
pública sobre diseño de impuestos. Supone curvas de indiferencia, restricción
presupuestaria y la idea de óptimo. Sirve también, sin cambiar nada, para
explicar por qué en los datos fiscales aparece gente amontonada justo debajo de
un umbral.

%% Secciones comunes a todos los manuales, edición española.
%% Se incluyen con \input{common-es}. Lo único que cambia entre applets es
%% \appletfolder, que es también el nombre de la carpeta y de la ruta en el sitio.

\section{Cómo abrirlo}

No hay que instalar nada. La aplicación entera es un solo archivo
\code{index.html}: no descarga librerías, no llama a ningún servidor y no
guarda nada fuera del navegador.

\subsection{Las tres maneras}

\begin{description}
\item[Doble clic sobre el archivo.] Es lo más rápido. Busque
  \code{\appletfolder/index.html} y ábralo. El navegador lo muestra desde el disco,
  con la barra de direcciones empezando por \code{file://}. Funciona sin
  conexión desde el primer momento.
\item[Desde la web.] Si alguien ya lo publicó, la dirección termina en
  \code{/\appletfolder/}. Una vez cargada la página, puede desconectarse: no
  necesita la red para nada más.
\item[Servido desde una carpeta.] Para una clase, poner la carpeta en
  cualquier servidor estático basta. Con Python instalado:
  \begin{quote}\code{python3 -m http.server 8000}\end{quote}
  y luego \code{http://localhost:8000/\appletfolder/} en el navegador.
\end{description}

\begin{amnote}
Guardar la página con \key{Ctrl}\,+\,\key{S} (o \key{Cmd}\,+\,\key{S} en
Mac) deja una copia propia que sigue funcionando sin conexión y que se puede
llevar en una memoria USB o repartir por correo.
\end{amnote}

\subsection{Idioma y tema}

Arriba a la derecha hay dos pares de botones.

\begin{ctrltable}{Control & Qué hace}
ES / EN & Cambia el idioma de toda la interfaz al vuelo. Nada se recalcula:
          puede cambiarlo en mitad de una explicación sin perder el estado. \\
Oscuro / Claro & Cambia el tema. El oscuro es el diseño original; el claro
          existe para proyectores y para imprimir. \\
\end{ctrltable}

Las dos elecciones se recuerdan en el navegador, así que la próxima vez se abre
como la dejó.

\subsection{Qué navegador}

Cualquiera que sea razonablemente reciente: Chrome, Edge, Firefox o Safari, en
ordenador, tableta o teléfono. En pantallas estrechas los paneles se apilan uno
sobre otro en lugar de ponerse al lado. No hace falta cuenta, ni permisos, ni
extensiones.


% ==========================================================================
\section{La pantalla}

\fig[1.0]{interface}{La pantalla completa: los controles a la izquierda, el
veredicto arriba, el plano consumo-ocio y el panel de renta neta abajo.}

\begin{description}
\item[Cómo se aplica.] Los dos esquemas, tipos marginales y tipos medios. Es el
  control que importa; todo lo demás es contexto.
\item[Tramos.] Tres tipos $t_0, t_1, t_2$ y dos umbrales $s_0, s_1$, todos
  sobre la renta del trabajo $wL$.
\item[Salario y renta.] El salario $w$, la renta no laboral $m_0$ y el precio
  del consumo $p_x$.
\item[Preferencias.] Cuatro familias y una casilla para escribir la suya.
\item[Eje vertical.] Si el eje del plano mide ocio o trabajo.
\item[Capas.] Qué se dibuja.
\end{description}

El veredicto de la banda superior nombra qué clase de óptimo ha salido:
tangencia, quiebre, borde del salto, no trabajar, o trabajar toda la dotación.

\subsection{Los dos paneles}

\begin{description}
\item[Consumo y ocio.] El mapa de utilidad bajo una rampa topográfica, lo
  inasequible velado, la frontera con sus quiebres o saltos, los umbrales, la
  curva de indiferencia por el óptimo, y el óptimo.
\item[Renta neta.] La renta neta frente a la bruta, con la recta de 45° de
  referencia. Aquí las preferencias no pintan nada, y ese es el punto: el salto
  es una propiedad del impuesto, no del trabajador. El conmutador lo cambia a
  \ui{Utilidad}, que dibuja $u(c(L), 1-L)$ a lo largo de la frontera.
\end{description}

\fig[0.62]{notch}{El panel de renta neta con tipos medios. En cada umbral la
renta neta cae en vertical: con los valores de partida, cruzar el primer umbral
la lleva de $20$ a $12$.}

\subsection{El tema oscuro}

\fig[0.95]{dark}{Los dos paneles en el tema oscuro, con tipos medios.}

% ==========================================================================
\section{El modelo}

\subsection{Los dos esquemas}

Sean $t_0, t_1, t_2$ los tipos y $s_0 \le s_1$ los umbrales, todos sobre la
renta del trabajo $Y = wL$.

\subsubsection{Tipos marginales}

Cada tipo grava solo la parte de la renta que cae dentro de su tramo:
\[
  T(Y) \;=\; t_0\min(Y, s_0)
        \;+\; t_1\max\bigl(0, \min(Y, s_1) - s_0\bigr)
        \;+\; t_2\max(0, Y - s_1).
\]
$T$ es continua y creciente, y la renta neta $Y - T(Y)$ es continua y cóncava.
La frontera presupuestaria es una poligonal: en cada umbral cambia de
pendiente, pero no se rompe.

\subsubsection{Tipos medios}

El tramo en el que se cae fija el tipo sobre \emph{toda} la renta:
\[
  T(Y) \;=\; Y \cdot t_{i(Y)},
  \qquad
  i(Y) = \begin{cases} 0 & Y < s_0\\ 1 & s_0 \le Y < s_1\\ 2 & Y \ge s_1\end{cases}
\]
Aquí $T$ salta en cada umbral, y la renta neta salta hacia abajo. Con los
valores de partida ($t_1 = 0{,}4$), cruzar $s_0 = 20$ lleva la renta neta de
$20$ a $12$: el trabajador que gana una unidad más se queda con ocho menos.

\subsection{Quiebre contra salto}

Los dos nombres se confunden con facilidad y la diferencia lo es todo.

\begin{ctrltable}{Nombre & Qué es}
Quiebre & La función es continua; la derivada salta. La frontera cambia de
  inclinación. Un óptimo puede caer exactamente ahí, y de hecho suele hacerlo. \\
Salto & La función misma es discontinua. La frontera se corta y hay un tramo de
  puntos que ningún $L$ alcanza. El óptimo se queda pegado justo por debajo. \\
\end{ctrltable}

\begin{amnote}
En la literatura fiscal, un quiebre se llama \emph{kink} y un salto
\emph{notch}. Los dos generan amontonamiento en los datos, pero por razones
distintas: en el quiebre porque la tasa marginal cambia; en el salto porque
cruzarlo empeora al contribuyente sin más.
\end{amnote}

\subsection{Cómo se encuentra el óptimo}

No con una condición de primer orden. En estas fronteras el mejor punto está
muy a menudo en un quiebre, donde la pendiente no existe y ninguna tangencia se
cumple; y con tipos medios está en el extremo derecho de una pieza, con un
punto estrictamente peor inmediatamente a su derecha.

Lo que se hace:

\begin{enumerate}
\item Se parte la frontera en piezas, una por tramo, con sus propios extremos.
\item Cada pieza se barre densamente y su mejor muestra se refina por sección
  áurea.
\item Se prueba \emph{cada vértice explícitamente}, y por los dos lados, cada
  uno con la fórmula de su propio tramo.
\item Gana el mejor de todos, y de dónde salió es lo que lo clasifica.
\end{enumerate}

\begin{amwatch}
Cada pieza lleva su propio tramo en lugar de deducirlo de la renta. Deducirlo
hace que $c(L)$ tenga un solo valor en un umbral, lo que pierde silenciosamente
el límite por la izquierda —que es el punto más alto de la frontera— y lo dejaba
fuera de la vista calculada.
\end{amwatch}

% ==========================================================================
\section{Los controles}

\subsection{Tramos, salario y renta}

\begin{ctrltable}{Deslizador & Qué es}
$t_0$, $t_1$, $t_2$ & Los tres tipos, de 0 a 0{,}95. \\
$s_0$, $s_1$ & Los dos umbrales, sobre la renta del trabajo. $s_1$ se mantiene
  igual o por encima de $s_0$: al revés invertiría el esquema. \\
$w$ & El salario. Fija dónde caen los umbrales en el eje del trabajo,
  $L = s/w$. \\
$m_0$ & Renta no laboral. Desplaza toda la frontera hacia arriba. \\
$p_x$ & Precio del consumo. Acotado lejos de cero. \\
\end{ctrltable}

\subsection{Preferencias}

\begin{ctrltable}{Familia & Forma}
Cobb--Douglas & $x^{\alpha} y^{1-\alpha}$ \\
CES & $(\alpha x^{\rho} + (1-\alpha) y^{\rho})^{1/\rho}$ \\
Cuasilineal & $x + 4\alpha\ln y$ \\
Personalizada & Lo que escriba \\
\end{ctrltable}

En la expresión $x$ es el consumo e $y$ el ocio. Con $\rho$ muy negativo la CES
se acerca a complementarios perfectos y aparecen óptimos en esquina.

\subsection{Eje vertical y capas}

El conmutador \ui{Eje vertical} cambia entre \ui{Ocio} y \ui{Trabajo}. La
utilidad se lee siempre sobre el ocio, así que voltear el eje relee de verdad
las preferencias en vez de reflejar un dibujo.

\fig[0.62]{labour-axis}{El plano con el eje vertical puesto en trabajo.}

\begin{ctrltable}{Capa & Qué dibuja}
Mapa de utilidad & El campo $u$ bajo una rampa de color. \\
Curvas de fondo & Curvas de nivel de la utilidad, independientes del óptimo. \\
Conjunto factible & Vela lo que no se puede pagar. \\
Curva por el óptimo & La curva de indiferencia que pasa por el óptimo. \\
Óptimo & El punto elegido. \\
Umbrales & Las verticales en $L = s_0/w$ y $L = s_1/w$. \\
Retícula & La cuadrícula de fondo. \\
\end{ctrltable}

% ==========================================================================
\section{El veredicto}

\begin{ctrltable}{Veredicto & Qué significa}
Óptimo interior & Tangencia: la relación marginal de sustitución iguala al
  salario neto del tramo. El trabajador está dentro de un tramo. \\
Óptimo en el quiebre & Cae exactamente sobre un umbral. A la izquierda todavía
  compensa trabajar más; a la derecha ya no. \\
Óptimo pegado al salto & Con tipos medios. La hora siguiente cruzaría el umbral
  y recaudaría el tipo más alto sobre toda la renta. \\
No trabajar es lo mejor & Esquina en $L = 0$: se vive de la renta no laboral. \\
Trabajar toda la dotación & Esquina en $L = 1$. \\
\end{ctrltable}

\fig[0.85]{kink-verdict}{El veredicto con los valores de partida y tipos
marginales: el óptimo cae sobre el primer umbral.}

% ==========================================================================
\section{Actividades para clase}

\begin{activity}{El mismo impuesto, dos fronteras}
\amset Valores de partida. Alterne entre \ui{Tipos marginales} y
\ui{Tipos medios} mirando el panel de la izquierda.
\amexp Con marginales la frontera se dobla en los umbrales. Con medios se
\emph{corta}: aparecen dos trozos de línea que no se tocan, y una franja
entera de consumos que ningún nivel de trabajo alcanza.
\par\smallskip
Los tipos, los umbrales, el salario y las preferencias son idénticos en los dos
casos. Lo único que cambió fue la frase «sobre la renta del tramo» por «sobre
toda la renta».
\end{activity}

\begin{activity}{Amontonarse justo debajo}
\amset Valores de partida, \ui{Tipos marginales}.
\amexp La lectura \ui{Trabajo L*} da $0{,}400$, que con $w = 50$ es exactamente
la renta bruta $20$: el primer umbral. El veredicto dice
«Óptimo en el quiebre».
\par\smallskip
Mueva $t_1$ arriba y abajo. El óptimo se queda clavado en el umbral para un
rango ancho de tipos, porque en un quiebre no hay una sola pendiente que
igualar sino un intervalo entero de ellas. Eso es lo que produce el
amontonamiento en los datos, y no es evasión: es el óptimo.
\end{activity}

\begin{activity}{Un salto que cuesta bienestar}
\amset Preferencias \ui{Cuasilineal}. Compare \ui{Tipos marginales} con
\ui{Tipos medios}.
\amexp Con marginales: $L^\ast = 0{,}800$, consumo $38{,}00$, renta bruta
$40{,}0$, impuesto $12{,}0$, utilidad $34{,}781$. El tipo medio es $30\%$ y el
marginal $80\%$.
\par\smallskip
Con medios, el mismo trabajador retrocede a $L^\ast = 0{,}400$, consumo
$30{,}00$, impuesto $0{,}0$, utilidad $28{,}978$. Recauda \emph{menos} impuesto
y el trabajador está \emph{peor}. Pregunte a quién beneficia ese diseño. A
nadie: es una pérdida pura, y es la razón por la que casi ningún sistema fiscal
usa tipos medios.
\end{activity}

\begin{activity}{Medir el salto}
\amset \ui{Tipos medios}, panel derecho en \ui{Impuesto}.
\amexp La lectura \ui{Mayor salto} da $-12{,}00$. Búsquelo en el dibujo: está
en el segundo umbral, $s_1 = 30$, donde el tipo pasa de $0{,}4$ a $0{,}8$ sobre
toda la renta. Justo debajo la renta neta es $30 - 12 = 18$; justo encima,
$30 - 24 = 6$.
\par\smallskip
El primer umbral también salta, de $20$ a $12$, pero menos. Baje $t_2$ hasta que
el mayor salto pase a ser el primero y vea cómo cambia la lectura.
\end{activity}

\begin{activity}{La utilidad se rompe en trozos}
\amset \ui{Tipos medios}, panel derecho en \ui{Utilidad}.
\amexp La curva $u(c(L), 1-L)$ aparece partida en arcos inconexos, uno por
tramo. Su punto más alto está en el extremo derecho de uno de ellos, no en un
máximo interior.
\par\smallskip
Es la misma información que el panel izquierdo, en una dimensión menos.
Pregunte dónde buscaría el máximo un estudiante que derivara e igualara a cero,
y por qué no lo encontraría.
\end{activity}

\begin{activity}{No trabajar}
\amset Preferencias \ui{CES}, $\rho = -6$.
\amexp El óptimo se va a $L^\ast = 0{,}000$: consumo $10{,}00$, que es
exactamente $m_0/p_x$. El veredicto dice «No trabajar es lo mejor».
\par\smallskip
Con $\rho$ muy negativo la CES se acerca a complementarios perfectos y el ocio
se vuelve casi imprescindible. Suba $w$ y mire cuánto hace falta para que
merezca la pena la primera hora.
\end{activity}

\begin{activity}{El eje no es la preferencia}
\amset Valores de partida. Cambie \ui{Eje vertical} de \ui{Ocio} a
\ui{Trabajo}.
\amexp El mapa de utilidad no es el reflejo del anterior. La utilidad se lee
siempre sobre el ocio, así que voltear el eje cambia de verdad lo que se está
dibujando y no solo cómo se mira.
\par\smallskip
Es una comprobación útil contra la costumbre de leer los dos ejes como si
fueran intercambiables.
\end{activity}

\begin{activity}{Diseñar un impuesto sin saltos}
\amset \ui{Tipos medios}. Objetivo: que la frontera no se corte.
\amexp Solo hay una manera: poner los tres tipos iguales. Cualquier diferencia
entre $t_0$, $t_1$ y $t_2$ produce un salto en el umbral correspondiente.
\par\smallskip
Con tipos marginales, en cambio, cualquier combinación creciente da una
frontera continua. Ahí está el argumento de diseño entero, y cabe en una frase.
\end{activity}

% ==========================================================================
\section{Sintaxis de expresiones}

En la casilla de utilidad, $x$ es el consumo e $y$ el ocio.

\begin{ctrltable}{Elemento & Qué se admite}
Operadores & \code{+ - * / \^{} ( )} \\
Funciones & \code{sqrt ln log log10 exp abs sin cos tan min max pow} \\
Constantes & \code{e}, \code{pi} \\
Multiplicación implícita & \code{2xy} es \code{2*x*y} \\
Asociatividad & \code{\^{}} liga por la derecha y más fuerte que el menos
  unario: \code{-x\^{}2} es $-(x^2)$ \\
\end{ctrltable}

Las derivadas parciales se calculan simbólicamente. Con \code{min}, \code{max} o
\code{abs} —justo los casos con quiebre— se pasa a diferencias centradas.

% ==========================================================================
\section{Diferencias con el original de GeoGebra}

\begin{description}
\item[Una sola base imponible.] \code{TAX\_LABOR} gravaba $wL$, pero la rama de
  tipos medios \code{l1} gravaba $wL + I_0$. Los dos esquemas no eran
  comparables justo en el momento en que uno quiere compararlos. Aquí los dos
  gravan la renta del trabajo.
\item[Las discontinuidades están modeladas.] El original escribía el impuesto
  como \code{If[...]} anidados, que GeoGebra dibuja con un segmento que salva el
  hueco. Aquí cada tramo es una pieza aparte, así que un salto sigue siendo un
  salto.
\item[Cada pieza lleva su tramo.] Deducir el tramo de la renta hace que $c(L)$
  tenga un solo valor en un umbral, lo que pierde el límite por la izquierda: el
  punto más alto de la frontera caía fuera de la vista calculada y se recortaba.
\item[Preferencias y un óptimo.] El original dibujaba solo el conjunto
  presupuestario. Sin un óptimo, el quiebre y el salto son formas y no
  predicciones, así que se añaden una función de utilidad, la curva de
  indiferencia por el óptimo y el veredicto sobre el amontonamiento.
\item[$p_X$ y $w$ empezaban en 0], dividiendo por cero en todas partes. Los dos
  están acotados lejos de él.
\item[$s_1 < s_0$] invertiría el esquema; el segundo umbral se mantiene igual o
  por encima del primero.
\item[Objetos muertos.] \code{s\_2} tenía deslizador y etiqueta pero no entraba
  en \code{TAX\_LABOR}; \code{a}, \code{b}, \code{f}, \code{ngreso}, \code{E},
  \code{E\_trace}, \code{Pendiente}, \code{B\_OCIO}, \code{B\_TRABAJO},
  \code{h\_TRABAJO} y el conjunto de puntos \code{Q}/\code{P} eran restos del
  applet del que se copió. No se trasladaron.
\end{description}

% ==========================================================================
\section{Problemas frecuentes}

\begin{ctrltable}{Síntoma & Qué pasa}
El plano sale en blanco & La expresión de utilidad no se entiende. El mensaje en
  rojo bajo la casilla dice si es sintaxis o variable. \\
Los umbrales no se ven & Están fuera del eje: con $w$ pequeño, $s/w$ puede pasar
  de 1. Suba $w$ o baje los umbrales. \\
Los dos esquemas dan lo mismo & Los tres tipos son iguales. Con tipos iguales no
  hay ni quiebre ni salto y los dos esquemas coinciden. \\
Mayor salto dice «---» & Está en tipos marginales, donde por construcción no hay
  saltos. \\
El óptimo no se mueve al cambiar un tipo & Está clavado en un quiebre. Es el
  comportamiento correcto y el asunto de la actividad 2. \\
\end{ctrltable}

%% Secciones comunes a todos los manuales, edición española.
%% Se incluyen con \input{common-es}. Lo único que cambia entre applets es
%% \appletfolder, que es también el nombre de la carpeta y de la ruta en el sitio.

\section{Publicarlo para una clase}

Sirve cualquier alojamiento estático: GitHub Pages, Netlify, un directorio web
de la universidad, o incluso una carpeta compartida servida por HTTP. Lo único
que hay que subir es el archivo \code{index.html} dentro de una carpeta con el
nombre que quiera que aparezca en la dirección.

El repositorio trae un flujo de trabajo de GitHub Actions
(\code{.github/workflows/pages.yml}) que publica todos los applets a la vez en
cada empuje a la rama principal. Para activarlo, en el repositorio:
\emph{Settings\,$\rightarrow$\,Pages\,$\rightarrow$\,Source: GitHub Actions}.
Es lo único que no puede hacer el propio flujo, porque crear un sitio de Pages
requiere permisos de administración que un \code{GITHUB\_TOKEN} nunca recibe.

\begin{amnote}
Para repartir el applet a estudiantes sin conexión fiable, mándeles el archivo
\code{index.html} directamente. Pesa menos que una fotografía y se abre con
doble clic.
\end{amnote}

\section{Licencia y procedencia}

El código de la aplicación es MIT. Puede usarlo, modificarlo y repartirlo, en
clase o fuera de ella, citando la procedencia.

Este applet es una reescritura autónoma de un applet original de GeoGebra. La
reescritura no es una traducción literal: donde el original tenía un error de
construcción, una división por cero alcanzable con los deslizadores, o un objeto
muerto heredado del applet del que se copió, esta versión lo corrige y lo
documenta. La sección \emph{Diferencias con el original de GeoGebra} enumera
esos cambios uno por uno.

Todo está escrito a mano y sin dependencias: el analizador de expresiones, la
derivación simbólica, el trazado de curvas de nivel y el dibujo. En particular
no se usa \code{eval}, de modo que lo que escribe un estudiante en una casilla
de texto no puede convertirse en código, y una política de seguridad de
contenidos estricta no rompe la página.


\end{document}
